\chapter{Simulation Model Details}
\label{app:B}

This appendix records the parameter values and structural details of the simulation model that
are summarised in \Cref{ch:methodology}.

\section{Workload multipliers}

Workload units for order \(i\) at stage \(s\) are computed by \Cref{eq:workload-units} from the
item count and the multipliers in \Cref{tab:multipliers}. The values differ by stage, which is
what produces stage-differentiated workload.

\begin{table}[htbp]
\centering
\caption[Workload multipliers by stage]{Workload multipliers. A multiplier of \num{1.0} is the
reference case. Urgency affects dispatch only.}
\label{tab:multipliers}
\small
\begin{tabular}{llrrr}
\toprule
\textbf{Attribute} & \textbf{Value} & \textbf{Picking} & \textbf{Packing} & \textbf{Dispatch} \\
\midrule
\multirow{3}{*}{Product family}
 & standard & 1.0 & 1.0 & 1.0 \\
 & fragile  & 1.1 & \textbf{1.8} & 1.1 \\
 & bulky    & \textbf{1.3} & 1.6 & 1.2 \\
\midrule
\multirow{3}{*}{Complexity}
 & low    & 0.9 & 0.8 & 0.9 \\
 & medium & 1.1 & 1.2 & 1.1 \\
 & high   & 1.4 & \textbf{1.7} & 1.4 \\
\midrule
\multirow{2}{*}{Urgency}
 & normal & 1.0 & 1.0 & 1.0 \\
 & urgent & 1.0 & 1.0 & \textbf{1.3} \\
\bottomrule
\end{tabular}
\end{table}

The entries in bold illustrate the point made in \Cref{sec:order-heterogeneity}: a fragile order
is only slightly harder to pick (\num{1.1}) but substantially harder to pack (\num{1.8}), while
a bulky order is relatively harder to pick than a fragile one. No single ordering of order
"difficulty" applies across the three stages.

\section{Service-time parameters}

\begin{table}[htbp]
\centering
\caption{Service-time parameters, per \Cref{eq:service-time}.}
\label{tab:service-params}
\small
\begin{tabular}{lrrr}
\toprule
\textbf{Parameter} & \textbf{Picking} & \textbf{Packing} & \textbf{Dispatch} \\
\midrule
Base time \(b_s\) (min)              & 0.50 & 0.50 & 0.50 \\
Minutes per workload unit \(c_s\)    & 0.90 & 0.70 & 0.65 \\
Minimum service time (min)           & 0.20 & 0.20 & 0.20 \\
Noise lower clip                     & 0.80 & 0.80 & 0.80 \\
Noise upper clip                     & 1.25 & 1.25 & 1.25 \\
\bottomrule
\end{tabular}
\end{table}

The multiplicative noise term is drawn from \(\mathcal{N}(1, 0.12)\) and clipped to
\([0.80, 1.25]\). For a reference order --- normal urgency, standard family, medium complexity,
three items --- these parameters produce expected service times of approximately
\SI{3.5}{\minute} at picking, \SI{2.0}{\minute} at packing and \SI{1.2}{\minute} at dispatch.

\section{Service-level and economic parameters}

\begin{table}[htbp]
\centering
\caption[Service-level and economic parameters]{Service-level thresholds, feasibility floors and
the economic parameters used in each of the two runs. The two runs use \emph{different}
economic assumptions; monetary values are never compared between them.}
\label{tab:sla-econ}
\small
\begin{tabularx}{\textwidth}{Xrr}
\toprule
\textbf{Parameter} & \textbf{Retrospective run} & \textbf{Forecast run} \\
\midrule
Urgent system-time SLA threshold (min) & 240 & 240 \\
Normal system-time SLA threshold (min) & \num{1440} & \num{1440} \\
Urgent attainment floor \(\theta_u\)  & 0.95 & 0.95 \\
Normal attainment floor \(\theta_n\)  & 0.80 & 0.80 \\
\midrule
Late urgent order penalty (EUR)       & \textbf{20} & \textbf{15} \\
Late normal order penalty (EUR)       & \textbf{5}  & \textbf{10} \\
Labour cost (EUR per hour)            & \textbf{15} & \textbf{18} \\
Paid hours per worker-month           & 160 & 160 \\
Monthly cost of one FTE (EUR)         & \num{2400} & \num{2880} \\
\midrule
Operating horizon \(H\) (min)         & \num{9600} & \num{9600} \\
Target utilisation \(\upsilon\)       & 0.85 & 0.85 \\
Candidates per month                  & 16 & 16 \\
Replications                          & 1 & 3 \\
\bottomrule
\end{tabularx}
\end{table}

\section{Structural assumptions}

\begin{itemize}
  \item Queues are unbounded; there is no blocking between stages.
  \item Workers are dedicated to a stage and cannot be reallocated during the month.
  \item Service is non-pre-emptive: an order in service is never interrupted.
  \item There is no travel time, setup between orders, batching or wave release.
  \item There is no worker learning, fatigue or absence.
  \item The simulation halts at \(H\); unfinished orders are backlog and count as service
        failures.
  \item Queue length is tracked by an event-driven accumulator giving an exact time-weighted
        average rather than a sampled approximation.
\end{itemize}

\section{Software environment}

\begin{table}[htbp]
\centering
\caption{Software environment.}
\label{tab:software}
\small
\begin{tabularx}{0.8\textwidth}{lX}
\toprule
\textbf{Component} & \textbf{Version / role} \\
\midrule
Python   & 3.12 \\
SimPy    & 4.1.1 --- process-based discrete-event simulation \\
NumPy    & 2.3.2 --- random number generation, numerical work \\
pandas   & 2.3.1 --- data handling and metric aggregation \\
PyTorch  & neural network implementation for RL-3; \textbf{version not recorded}
           --- see the reproducibility limitation in \Cref{app:E} \\
Matplotlib & 3.10.7 --- figure generation \\
\bottomrule
\end{tabularx}
\end{table}
